% Documento autocontenido. Compilar dos veces con pdfLaTeX.
\documentclass[11pt,a4paper]{article}
\usepackage[utf8]{inputenc}
\usepackage[T1]{fontenc}
\usepackage{mathpazo}
\usepackage{amsmath,amssymb}
\usepackage[a4paper,margin=20mm]{geometry}
\usepackage[expansion=false]{microtype}
\usepackage[hidelinks]{hyperref}
\hypersetup{pdftitle={14 temas de tesis de licenciatura en información cuántica},pdfauthor={Propuesta para Marc Urcia y Danesh Morales}}
\pagestyle{empty}
\setlength{\parindent}{0pt}
\setlength{\parskip}{4pt}
\newcommand{\tema}[5]{%
\noindent\begin{minipage}{\linewidth}\raggedright
\textbf{#1. #2}\par
\textbf{Pregunta:} #3\par
#4\par
\textbf{Producto:} #5
\end{minipage}\par\vspace{9pt}}
\begin{document}
\raggedright

\tema{1}{De códigos clásicos a códigos cuánticos: construcción CSS}
{¿Cómo se transforma un código corrector clásico en uno que protege información cuántica?}
{El estudiante desarrollaría códigos binarios, matrices de chequeo, estabilizadores y la construcción CSS, usando Steane como ejemplo central. Implementaría un programa que verifique conmutación, número de qubits lógicos y distancia en ejemplos pequeños. Como extensión, estudiaría un conjunto acotado de modificaciones de las matrices de chequeo~\cite{css}.}
{Una exposición matemática rigurosa acompañada de una herramienta reproducible para analizar códigos. Es una buena entrada para estudiantes de \textbf{Matemática con álgebra lineal y algo de álgebra abstracta}. El curso de IBM proporciona un punto de partida~\cite{ibmqec}.}

\tema{2}{Decodificación adaptada al ruido}
{¿Cuánto ayuda conocer las probabilidades de los errores para elegir una recuperación?}
{El estudiante trabajaría con el código $[[5,1,3]]$ y ruido de Pauli. Enumeraría los errores y compararía recuperación de peso mínimo con máxima probabilidad de la clase lógica, sumando errores equivalentes por estabilizadores. Después variaría el sesgo del ruido y la precisión con que se conocen sus probabilidades~\cite{adapted}.}
{Un mapa de rendimiento y sensibilidad del decodificador al modelo de ruido. Adecuado para estudiantes de \textbf{Matemática o Física con probabilidad, álgebra lineal y programación}, interesados en algoritmos y decisiones basadas en datos.}

\tema{3}{Topología y protección de información: el código tórico}
{¿Cómo se relacionan los ciclos de una superficie con los errores lógicos de un código cuántico?}
{El estudiante construiría el código tórico mediante vértices, aristas y caras. Explicaría la conmutación de los estabilizadores y la relación entre ciclos no triviales y operadores lógicos. Programaría ejemplos pequeños con matrices binarias y visualizaciones de cadenas de errores, sin almacenar el estado cuántico completo~\cite{toric}.}
{Un desarrollo matemático del código acompañado de una herramienta visual que verifique sus propiedades. Atractivo para estudiantes de \textbf{Matemática con álgebra lineal e interés en topología y teoría de grafos}; la homología elemental se aprendería durante el proyecto.}

\tema{4}{Tomografía de un qubit y modelos de ruido}
{¿Cuándo basta un modelo de errores de Pauli para describir un proceso cuántico?}
{El estudiante reconstruiría un canal de un qubit a partir de preparaciones y mediciones en distintas bases. Compararía un modelo de Pauli con un ajuste general físicamente válido, primero con datos simulados y luego, opcionalmente, en IBM. Evaluaría predicciones en datos reservados y el efecto de errores de preparación y lectura~\cite{tomography}.}
{Una comparación de precisión, costo de medición y límites del modelo. Apropiado para estudiantes de \textbf{Matemática aplicada o Física con álgebra lineal, estadística y métodos numéricos}, interesados en problemas inversos y optimización convexa.}

\tema{5}{Benchmarking y predicción del rendimiento}
{¿Pueden dos modelos con infidelidad media similar producir resultados distintos en un circuito?}
{El estudiante implementaría randomized benchmarking de un qubit y compararía sobrerrotaciones coherentes con ruido estocástico. Ajustaría los decaimientos, estimaría su incertidumbre y estudiaría una familia pequeña de circuitos repetidos para determinar qué información aporta la métrica de benchmarking sobre su rendimiento~\cite{rb}.}
{Ejemplos cuantitativos de la capacidad predictiva y los límites de una métrica promedio. Adecuado para estudiantes de \textbf{Física con estadística, métodos numéricos y mecánica cuántica básica}, interesados en caracterización experimental y análisis de datos.}

\tema{6}{Errores de lectura correlacionados y mitigación}
{¿Cuándo es insuficiente calibrar la lectura de cada qubit por separado?}
{El estudiante trabajaría con dos o tres qubits y compararía datos sin mitigación, un modelo de lectura independiente y una matriz conjunta. Variaría las correlaciones y el número de mediciones, contabilizando tanto calibración como validación. Una ejecución pequeña en IBM permitiría contrastar los modelos con datos reales~\cite{readout}.}
{Criterios para elegir un modelo de lectura según precisión y costo experimental. Es una entrada accesible para estudiantes de \textbf{Matemática o Física con probabilidad, álgebra lineal y Python}, interesados en calibración, incertidumbre y análisis numérico.}

\tema{7}{Código de repetición en IBM: frecuencia de recuperación}
{¿Cuándo compensa corregir con mayor frecuencia el ruido añadido por la recuperación?}
{El estudiante implementaría un código de repetición con tres qubits de datos y dos auxiliares. Compararía decodificación final y recuperación durante el circuito a duración total fija, primero en simulación y después en IBM~\cite{wootton,ibmrep}. Estudiaría errores de tipo $X$ y distinguiría conservación en la base computacional de conservación de coherencia.}
{Curvas de rendimiento que muestren cuándo ayuda la recuperación. Adecuado para estudiantes de \textbf{Física o Matemática aplicada con programación, álgebra lineal y fundamentos de circuitos cuánticos}, interesados en experimentación remota.}

\tema{8}{Compilación aleatorizada y acumulación del error}
{¿Cómo cambia la acumulación del error al aleatorizar la implementación de un circuito?}
{El estudiante reproduciría randomized compiling en circuitos pequeños con sobrerrotaciones controladas. Verificaría que las operaciones aleatorias y sus compensaciones preservan el circuito ideal. Compararía rendimiento y variabilidad con y sin aleatorización, estudiando cuántas realizaciones hacen falta a presupuesto total de mediciones fijo~\cite{rc}.}
{Una implementación reproducible y curvas de precisión frente a costo. Apropiado para estudiantes de \textbf{Matemática o Física con álgebra lineal, probabilidad y nociones de grupos}, interesados en algoritmos aleatorios y control del error.}

\tema{9}{Memoria del entorno y eficacia de QEC}
{¿Cómo cambia la recuperación cuando el entorno conserva información entre dos intervalos?}
{El estudiante simularía un código pequeño acoplado a un qubit de entorno. Compararía dos intervalos con el mismo entorno frente a un control donde se reinicializa, manteniendo la interacción elemental. Mediría fidelidad lógica y estadísticas de síndromes. El trabajo de Oreshkov y Brun sobre corrección continua serviría de antecedente conceptual~\cite{memory}.}
{Un estudio delimitado del efecto de la memoria sobre la recuperación. Recomendado para estudiantes de \textbf{Física o Matemática aplicada con mecánica cuántica, álgebra lineal sólida y programación científica}, interesados en sistemas abiertos y dinámica.}

\tema{10}{Desacoplamiento dinámico durante tiempos de espera}
{¿Cuándo ayudan las secuencias de control a preservar la coherencia de un qubit?}
{El estudiante compararía espera libre, eco y XY4 bajo desfase cuasiestático o temporalmente correlacionado, manteniendo igual duración total. Incluiría errores de las operaciones de control y estudiaría distintos estados iniciales. Una campaña pequeña en IBM podría complementar las simulaciones~\cite{dd}.}
{Un mapa de los regímenes donde las secuencias ayudan o perjudican. Adecuado para estudiantes de \textbf{Física con mecánica cuántica, ecuaciones diferenciales y programación}, interesados en teoría de control y fenómenos de coherencia.}

\tema{11}{Medida de Haar y operaciones cuánticas aleatorias}
{¿Qué significa elegir uniformemente una operación cuántica y cómo se genera esa distribución?}
{El estudiante introduciría la invariancia de la medida de Haar y generaría unitarias de dimensiones dos y cuatro mediante matrices gaussianas y factorización QR con corrección de fases. Compararía este procedimiento con un muestreo ingenuo de ángulos, usando momentos y pruebas numéricas de invariancia~\cite{haar}.}
{Una justificación matemática del algoritmo y una comparación reproducible entre muestreo correcto y sesgado. Atractivo para estudiantes de \textbf{Matemática con álgebra lineal, probabilidad e interés en teoría de la medida y grupos}, sin exigir formación cuántica avanzada.}

\tema{12}{Diseños unitarios y estimación de fidelidad}
{¿Qué promedios de Haar puede reproducir un conjunto finito de operaciones?}
{El estudiante estudiaría diseños unitarios de órdenes uno y dos, comparando Pauli y Clifford para un qubit. Analizaría sus momentos y aplicaría los estados generados mediante Clifford a la estimación de fidelidad promedio de un canal. Distinguiría la igualdad exacta de momentos de la incertidumbre debida al muestreo finito~\cite{designs}.}
{Una exposición con demostraciones y una comparación de estimadores. Apropiado para estudiantes de \textbf{Matemática con álgebra lineal, probabilidad y algo de álgebra abstracta}, interesados en sustituir integrales por conjuntos finitos y conectar simetrías con caracterización.}

\tema{13}{Simetrías y subespacios libres de decoherencia}
{¿Cómo protege una simetría la información y qué ocurre cuando se rompe ligeramente?}
{El estudiante analizaría dos qubits sujetos a desfase colectivo y demostraría la protección del subespacio generado por $|01\rangle$ y $|10\rangle$. Introduciría una pequeña desigualdad entre sus acoplamientos y estudiaría la pérdida de fidelidad, delimitando qué errores respeta la protección ideal~\cite{dfs}.}
{Una derivación de la protección por simetría y un estudio de su robustez. Atractivo para estudiantes de \textbf{Física o Matemática con álgebra lineal y mecánica cuántica básica}, interesados en subespacios invariantes, representaciones y sistemas abiertos.}

\tema{14}{Estados aleatorios y entrelazamiento bipartito}
{¿Cómo depende el entrelazamiento típico de un estado puro aleatorio de las dimensiones de sus subsistemas?}
{El estudiante generaría estados puros de Haar, calcularía su descomposición de Schmidt y estudiaría la pureza de los estados reducidos. Derivaría la pureza media y la contrastaría con histogramas y promedios numéricos para dimensiones pequeñas. Como extensión, podría analizar la entropía de von Neumann~\cite{randomstates}.}
{Una derivación analítica acompañada de simulaciones reproducibles. Adecuado para estudiantes de \textbf{Matemática o Física con álgebra lineal y probabilidad}, interesados en matrices aleatorias, teoría de la información y propiedades típicas de sistemas cuánticos.}

\clearpage
\renewcommand{\refname}{Referencias}
\begingroup
\small
\begin{thebibliography}{99}
\setlength{\itemsep}{5pt}
\setlength{\parskip}{0pt}
\interlinepenalty=10000
\bibitem{css} A. R. Calderbank y P. W. Shor (1996). ``Good quantum error-correcting codes exist''. \emph{Physical Review A} \textbf{54}, 1098. \href{https://doi.org/10.1103/PhysRevA.54.1098}{DOI}; \href{https://arxiv.org/abs/quant-ph/9512032}{arXiv}.
\bibitem{ibmqec} IBM Quantum Learning. \emph{Foundations of quantum error correction}. \href{https://learning.quantum.ibm.com/course/foundations-of-quantum-error-correction}{Curso: estabilizadores, CSS y código tórico}.
\bibitem{adapted} A. S. Fletcher, P. W. Shor y M. Z. Win (2008). ``Structured near-optimal channel-adapted quantum error correction''. \emph{Physical Review A} \textbf{77}, 012320. \href{https://doi.org/10.1103/PhysRevA.77.012320}{DOI}; \href{https://arxiv.org/abs/0708.3658}{arXiv}.
\bibitem{toric} A. Yu. Kitaev (2003). ``Fault-tolerant quantum computation by anyons''. \emph{Annals of Physics} \textbf{303}, 2--30. \href{https://doi.org/10.1016/S0003-4916(02)00018-0}{DOI}; \href{https://arxiv.org/abs/quant-ph/9707021}{arXiv}.
\bibitem{tomography} I. L. Chuang y M. A. Nielsen (1997). ``Prescription for experimental determination of the dynamics of a quantum black box''. \emph{Journal of Modern Optics} \textbf{44}, 2455--2467. \href{https://doi.org/10.1080/09500349708231894}{DOI}; \href{https://arxiv.org/abs/quant-ph/9610001}{arXiv}.
\bibitem{rb} E. Magesan, J. M. Gambetta y J. Emerson (2011). ``Scalable and Robust Randomized Benchmarking of Quantum Processes''. \emph{Physical Review Letters} \textbf{106}, 180504. \href{https://doi.org/10.1103/PhysRevLett.106.180504}{DOI}; \href{https://arxiv.org/abs/1009.3639}{arXiv}.
\bibitem{readout} F. B. Maciejewski, Z. Zimborás y M. Oszmaniec (2020). ``Mitigation of readout noise in near-term quantum devices by classical post-processing based on detector tomography''. \emph{Quantum} \textbf{4}, 257. \href{https://doi.org/10.22331/q-2020-04-24-257}{DOI}; \href{https://arxiv.org/abs/1907.08518}{arXiv}.
\bibitem{wootton} J. R. Wootton y D. Loss (2018). ``Repetition code of 15 qubits''. \emph{Physical Review A} \textbf{97}, 052313. \href{https://doi.org/10.1103/PhysRevA.97.052313}{DOI}; \href{https://arxiv.org/abs/1709.00990}{arXiv}.
\bibitem{ibmrep} IBM Quantum Documentation. \emph{Repetition codes}. \href{https://quantum.cloud.ibm.com/docs/en/tutorials/repetition-codes}{Tutorial de circuitos dinámicos}. Recursos web consultados el 1 de noviembre de 2026.
\bibitem{rc} J. J. Wallman y J. Emerson (2016). ``Noise tailoring for scalable quantum computation via randomized compiling''. \emph{Physical Review A} \textbf{94}, 052325. \href{https://doi.org/10.1103/PhysRevA.94.052325}{DOI}; \href{https://arxiv.org/abs/1512.01098}{arXiv}.
\bibitem{memory} O. Oreshkov y T. A. Brun (2007). ``Continuous quantum error correction for non-Markovian decoherence''. \emph{Physical Review A} \textbf{76}, 022318. \href{https://doi.org/10.1103/PhysRevA.76.022318}{DOI}; \href{https://arxiv.org/abs/0705.2342}{arXiv}.
\bibitem{dd} L. Viola y S. Lloyd (1998). ``Dynamical suppression of decoherence in two-state quantum systems''. \emph{Physical Review A} \textbf{58}, 2733--2744. \href{https://doi.org/10.1103/PhysRevA.58.2733}{DOI}; \href{https://arxiv.org/abs/quant-ph/9803057}{arXiv}.
\bibitem{haar} F. Mezzadri (2007). ``How to generate random matrices from the classical compact groups''. \emph{Notices of the AMS} \textbf{54}, 592--604. \href{https://arxiv.org/abs/math-ph/0609050}{arXiv}.
\bibitem{designs} C. Dankert, R. Cleve, J. Emerson y E. Livine (2009). ``Exact and approximate unitary 2-designs and their application to fidelity estimation''. \emph{Physical Review A} \textbf{80}, 012304. \href{https://doi.org/10.1103/PhysRevA.80.012304}{DOI}; \href{https://arxiv.org/abs/quant-ph/0606161}{arXiv}.
\bibitem{dfs} D. A. Lidar (2012). \emph{Review of Decoherence Free Subspaces, Noiseless Subsystems, and Dynamical Decoupling}. \href{https://arxiv.org/abs/1208.5791}{arXiv:1208.5791}.
\bibitem{randomstates} O. Giraud (2006). \emph{Distribution of bipartite entanglement for random pure states}. \href{https://arxiv.org/abs/quant-ph/0611285}{arXiv:quant-ph/0611285}.
\end{thebibliography}
\endgroup
\end{document}
